\documentclass[../../../include/open-logic-section]{subfiles}

\begin{document}

\olfileid[tr]{sth}{z}{milestone}
\olsection{$\Z^-$: Bir Dönüm Noktası}

\stagesinf{} ilkesini sonraki kesimde yeniden ele alacağız. Ancak Sonsuzluk
Aksiyomu'yla önemli bir dönüm noktasına ulaştık. Artık $\Zminus$ teorisi
için gerekli bütün aksiyomlara sahibiz. Ayrıntılı olarak:

\begin{defn}
$\Zminus$ teorisinin aksiyomları Kaplamsallık, Birleşim, Çiftler, Kuvvet
Kümeleri, Sonsuzluk ve Ayırma şemasının bütün örnekleridir.
\end{defn}

Ad, \emph{Zermelo} küme teorisini (daha sonra karşılaşacağımız bir şey
\emph{çıkarılmış} olarak) belirtir. Zermelo bu onuru hak eder; çünkü söz
konusu teoriyi esas olarak \citeyear{Zermelo1908Untersuchungen} yılında
formüle etmiştir.\footnote{Tarih ve teknik ayrıntılar üzerine ilginç
yorumlar için \citet[Ek A]{Potter2004} kaynağına bakın.}

Bu teori, muazzam miktarda matematik yapmamıza yetecek kadar güçlüdür.
Özellikle \olref[sfr][][]{part} boyunca geriye dönüp naif biçimde
yaptığımız her şeyin $\Zminus$ içinde daha biçimsel olarak yapılabileceğine
kendinizi \emph{ikna etmelisiniz}. (Bunu bir süre yaptıktan sonra ileri
atlayıp \olref[sth][z][arbintersections]{sec} kesimini okumak
isteyebilirsiniz.) Bundan böyle, ayrıca belirtmeden, en az $\Zminus$
içinde çalıştığımızı kabul edeceğiz.

\end{document}
